% T16 source: Практический пример.tex
\subsection{Illustrative Planning Scenario and Reproduction by EXP-00}
\label{sec:practical-example}

Consider a synthetic illustrative scenario involving the development of a small
request-processing service. The fixed workload comprises an API contract, a
persistence layer, business logic, background processing, observability, tests,
and packaging. On a notional scale, it equals $38$ story points, or $19$ units
of $Z$; the normalization $X=4$ hours yields a sequential no-AI baseline time of
$T_h=76$ hours. This relationship is used solely to make the scenario
transparent and is not an empirical conversion from story points to hours.

The task-DAG contains the sequential core
$1\to2\to3\to4\to6$ and two side branches: observability with packaging, and
work that becomes available once the interfaces have been fixed. Four variants
of work organization are compared. In Variants 1--3, the operating mode and
configured parallelism change, and the coefficients $k_i,h_i$ therefore change
as well. Variant 4 retains the configuration of Variant 3 but decomposes the
monolithic testing task into three parts with exact dependencies. Only the
comparison of Variants 3 and 4 isolates the effect of decomposition with all
other factors held constant.

\begin{table}[ht]
\centering
\caption{Condensed comparison of the illustrative-scenario variants}
\label{tab:compact-practical}
\begin{tabular}{lcccc}
\toprule
Metric & Variant 1 & Variant 2 & Variant 3 & Variant 4 \\
\midrule
$P$ & 1 & 2 & 4 & 4 \\
Mode & interactive & mixed & delegated & delegated \\
Decomposition & coarse & coarse & coarse & split tests \\
$W_4$, h & 65.0 & 59.8 & 53.6 & 53.6 \\
$W_4/P$, h & 65.0 & 29.9 & 13.4 & 13.4 \\
$L_4$, h & 51.8 & 47.8 & 43.2 & 36.7 \\
$H$, h & 65.0 & 38.4 & 19.0 & 19.0 \\
$B_4$, h & 65.0 & 47.8 & 43.2 & 36.7 \\
$T^*$, h & 65.0 & 47.8 & 43.2 & 36.7 \\
Speedup & $1.17\times$ & $1.59\times$ & $1.76\times$ & $2.07\times$ \\
\bottomrule
\end{tabular}
\end{table}

In EXP-00, the exact solver returned status \texttt{OPTIMAL} for Variants 2--4.
In each case, a feasible schedule with $T(\pi)=B_4$ was exhibited; hence, the
chain $B_4\le T^*\le T(\pi)$ proves $T^*=B_4$ specifically for these scenarios.
The equality is not assumed to be a universal property of M4.

\subsubsection{Local Blocking and the Human Resource}

In Variant 2, two agent slots are served by a single developer. A naive
schedule that is feasible with respect to the DAG and the number of agents
contains overlapping human-attention phases and is therefore infeasible under
M4. The corrected schedule in Figure~\ref{fig:practical-gantt} sequences the
task-specification and review phases. Task 7 waits $22.2$ hours for review and
retains the second slot, but this wait proceeds in parallel with the critical
branch and does not increase the project makespan. This example shows why
$Q(\pi)$ alone does not determine the makespan: the placement of the wait in the
resource schedule matters.
\begin{figure}[ht]
\centering
\includegraphics[width=\textwidth]{candidate_variant2_schedule_comparison.png}
\caption{Variant 2 before and after explicitly accounting for the human
resource. On the left, the original schedule contains overlapping
human-attention phases and is infeasible under M4. On the right, the developer's
service is serialized; the waiting task retains its agent slot but does not
increase the overall makespan.}
\label{fig:practical-gantt}
\end{figure}

\FloatBarrier

\subsubsection{Critical-Path Decomposition}

In Variant 3, the entire testing task waits for completion of the background
worker. In Variant 4, unit and integration tests begin once their respective
predecessors have completed, while only retry testing remains at the end of the
critical path. With $A=34.6$, $H=19.0$, and $W_4=53.6$ unchanged, the critical
path and exact makespan decrease from $43.2$ to $36.7$ hours. The effect is
$6.5$ hours, or approximately $15\%$ of the Variant 3 makespan; speedup relative
to the no-AI baseline increases from $1.76\times$ to $2.07\times$.

\begin{figure}[ht]
\centering
\includegraphics[width=\textwidth]{exp00_dag_variants_3_4.png}
\caption{Task-DAG and critical path before and after refining the dependencies
of the testing task. Total work is preserved, while unnecessary serialization
is reduced.}
\label{fig:practical-dag}
\end{figure}

\FloatBarrier

Complete task parameters, phase intervals, the human-resource check, and
elementary M0--M1 examples are provided in Appendix~\ref{app:examples}.
